\section{Introduction}
\label{sec:introduction}

Large Language Models (LLMs) deployed as autonomous agents---with tool access, persistent memory, and multi-turn interaction---create an attack surface that no existing defense adequately addresses. The fundamental challenges span five categories:

\begin{enumerate}[leftmargin=*]
\item \textbf{Prompt injection remains unsolved.} No production system reliably prevents instruction override via adversarial input~\cite{greshake2023}.
\item \textbf{Agentic attacks compound.} Tool chains create exponential attack surface: $N$ tools yield $O(N!)$ possible action chains.
\item \textbf{Model integrity is unverifiable.} Goldwasser and Kim~\cite{goldwasser2022} prove that backdoor detection is mathematically impossible under standard cryptographic assumptions.
\item \textbf{Semantic identity defeats classification.} Malicious and benign intent can produce identical text (``How do I mix bleach and ammonia?'').
\item \textbf{Defense layers conflict.} Provenance tracking and semantic transduction are architecturally incompatible without novel primitives.
\end{enumerate}

Existing commercial solutions---Lakera Guard, Meta Prompt Guard~\cite{promptguard2024}, NVIDIA NeMo Guardrails~\cite{nemo2023}, and LLM Guard---are all variants of content-level filtering: classify input as safe or dangerous, then block or allow. This approach has a fundamental ceiling, bounded by the impossibility results above.

\paragraph{Our approach.} Rather than incrementally improving content classification, we apply a methodology inspired by recent breakthroughs in AI-driven scientific discovery. AlphaFold~\cite{alphafold2021} did not invent new physics---it systematically explored 170,000 known protein structures and found patterns that 500,000 biologists missed due to cognitive limitations. Similarly, we systematically analyze \textbf{58 security paradigms from 19 scientific domains} and identify cross-domain combinations that have never been applied to LLM security.

This cross-domain synthesis yields seven novel security primitives (Table~\ref{tab:primitives}), five of which are genuinely new inventions with zero prior art, confirmed via 51 independent searches across GitHub and academic databases.

\begin{table}[t]
\centering
\small
\caption{The seven security primitives introduced in this work.}
\label{tab:primitives}
\begin{tabular}{@{}cp{3.0cm}cc@{}}
\toprule
\# & Primitive & Acr. & Nov. \\
\midrule
1 & Provenance-Annotated Sem.\ Reduction & \pasr & New \\
2 & Capability-Attenuating Flow Labels & \cafl & New \\
3 & Goal Predictability Score & \gps & New \\
4 & Adversarial Argumentation Safety & \aas & New \\
5 & Intent Revelation Mechanisms & \irm & New \\
6 & Model-Irrelevance Containment & \mire & New \\
7 & Temporal Safety Automata & \tsa & Adapted \\
\bottomrule
\end{tabular}
\end{table}

\paragraph{Contributions.} This paper makes the following contributions:
\begin{itemize}[leftmargin=*]
\item A systematic methodology for cross-domain paradigm synthesis applied to LLM security, analyzing 58 paradigms from 19 domains (\S\ref{sec:background}).
\item Seven novel security primitives, with formal definitions and architectural integration (\S\ref{sec:pasr}--\S\ref{sec:mire}).
\item A paradigm shift from backdoor \emph{detection} (proven impossible) to backdoor \emph{containment} via \mire (\S\ref{sec:mire}).
\item Evaluation against 250,000 simulated attacks across 15 categories, achieving approximately 98.5\% containment (\S\ref{sec:evaluation}).
\item Confirmation of zero prior art via exhaustive search across 51 cross-domain intersections (\S\ref{sec:related}).
\end{itemize}
